\documentclass{article}
\usepackage[utf8]{inputenc}
\usepackage[italian]{babel}
\usepackage{xcolor}
\usepackage[ruled,vlined,linesnumbered]{algorithm2e}

% Definizione dei colori per evidenziare le sezioni
\definecolor{basecolor}{RGB}{180, 0, 0}      % Rosso per il Caso Base
\definecolor{bodycolor}{RGB}{0, 100, 180}    % Blu per il Corpo
\definecolor{recurcolor}{RGB}{0, 130, 50}    % Verde per la Ricorsione

\begin{document}

\title{Algoritmi ricorsivi}
\author{Riccardo Di Michele}
\maketitle

\section*{Costo di un algoritmo ricorsivo}
Partiamo dalle basi, visto che la professoressa ci tiene. \\
Un \textbf{algoritmo} è un insieme di passaggi elementari e non ambigui 
che hanno come scopo la risoluzione di
un problema algoritmico, ovvero un problema risolvibili in un numero
finito di passaggi.
Un \textbf{algoritmo ricorsivo} è invece un algoritmo che 
\textit{richiama se stesso}, più nello specifico, la struttura di questo
tipo di algoritmo presenta: \textbf{caso base}, ovvero la condizione che
ci fa "uscire" dal ciclo di chiamate ricorsive, il \textbf{corpo}, gli
effettivi passagi che l'algoritmo esegue ad ogni iterazione, e le 
\textbf{chiamate} alla stessa funzione, con input 
di \textit{inferiore dimensione}.

\vspace{0.3cm}

\begin{algorithm}[H]
\caption{Search the value $v$ among the $n$ ordered elements of $S$}
\SetKwFunction{FBinarySearch}{BINARYSEARCH}

\SetKwProg{Fn}{function}{}{}
\Fn{\FBinarySearch{$S, v, first, last$}}{
    
    % --- CASO BASE 1 ---
    \textcolor{basecolor}{\If{$first > last$}{
        \Return $0$ \tcp*{Caso Base: elemento non trovato}
    }}
    
    % --- CORPO DELL'ALGORITMO ---
    \textcolor{bodycolor}{$m \leftarrow \lfloor \frac{first + last}{2} \rfloor$\; \tcp*{Corpo: calcolo indice medio}}
    
    % --- CASO BASE 2 ---
    \textcolor{basecolor}{\If{$S[m] = v$}{
        \Return $m$ \tcp*{Caso Base: elemento trovato}
    }}
    
    \Else{
        \If{$S[m] > v$}{
            % --- CHIAMATA RICORSIVA 1 ---
            \textcolor{recurcolor}{\Return \FBinarySearch{$S, v, first, m - 1$} \\ \tcp*{Chiamata Ricorsiva (sinistra)}}
        }
        \Else{
            % --- CHIAMATA RICORSIVA 2 ---
            \textcolor{recurcolor}{\Return \FBinarySearch{$S, v, m + 1, last$} \\ \tcp*{Chiamata Ricorsiva (destra)}}
        }
    }
}
\end{algorithm}

\vspace{0.3cm}

\noindent Prendendo come esempio l'algoritmo di ricerca binaria possiamo notare
tutti e tre i casi analizzati prima. Ma il cuore della materia è sempre quello,
quindi anche nel caso di quelli ricorsivi ci domandiamo: 
\textbf{quanto costa questo algoritmo?} \\
Negli algoritmi ricorsivi il costo si calcola con la seguente 
\textbf{relazione di ricorrenza}:
\[ T(n) = f(n) + k \cdot T(n'_k) \]

\noindent Dove $f(n)$ è il costo del \textbf{corpo}, mentre 
$k \cdot T(n'_k)$ sono il numero di \textbf{chiamate} effettuate per 
iterazione ($k$), moltiplicato per un $n$ di dimensione inferiore. \\
Nel caso del \textit{binary search} il corpo ha un costo di $O(1)$, in 
quanto ha solo una operazione di assegnazione, e il costo delle chiamate
è invece pari a $T(n/2)$; il valore viene dalla stessa logica discussa nel 
capitolo su questo algoritmo. \\

\[ T(n) = O(1) + T(n/2) \]

Notare come le chiamate ricorsive all'interno del codice sono 
\textit{tecnicamente} due ma nella formula il coefficiente $k$ è 
pari a $1$, questo perchè contiamo solo le chiamate ricorsive effettuate ad 
\textit{ogni iterazione}. Poichè abbiamo un solo $if$ soltanto una 
delle due chiamate verrà eseguita; se per esempio non avessimo avuto alcun
$if$ allora avremmo scritto un costo tale a $2T(n/2)$.

\section*{Metodi di scomposizione}
Adesso che abbiamo la nostra relazione di ricorrenza cosa ci facciamo? \\
La relazione così non ci dice nulla, ci serve esprimerla \textit{in funzione}
dell'input. Ci aiuteremo quindi usando dei 
\textbf{\textit{metodi di scomposizione}}.



\end{document}